\section{Lo Spettro Elettromagnetico}

%=======================================================================

\begin{frame}{Spettro elettromagnetico}

    Tutte le onde EM condividono la stessa natura, differiscono solo per $\lambda$ e $f$
    
    \medskip
    \centering
    
    \renewcommand{\arraystretch}{1.3}
    \arrayrulecolor{black}
    
    \begin{tabular}{|l|l|l|}
    \hline
    \rowcolor{maincolor!20}
    \textbf{Radiazione} & \textbf{$\lambda$ tipica} & \textbf{Applicazione} \\
    \hline
    Onde radio & da cm a km & Telecomunicazioni \\
    \hline
    Microonde & mm ÷ cm & Radar, Wi-Fi, forni \\
    \hline
    Infrarosso & \SI{780}{nm} ÷ \SI{1}{mm} & Termografia, telecomandi \\
    \hline
    \rowcolor{statalegreen!20}
    \textbf{Luce visibile} & \textbf{380 ÷ 780 nm} & \textbf{Visione} \\
    \hline
    Ultravioletto & 10 ÷ 380 nm & Sterilizzazione, abbronzatura \\
    \hline
    Raggi X & 0.01 ÷ 10 nm & Diagnostica medica \\
    \hline
    Raggi Gamma &  $<$ 0.01 nm & Fisica nucleare \\
    \hline
    \end{tabular}

\end{frame}

%=======================================================================

\begin{frame}{La percezione del Colore}
    \begin{itemize}
        \item La luce bianca contiene \textbf{tutte le lunghezze d'onda} del visibile
        \item Gli oggetti \textbf{assorbono} alcune $\lambda$ e \textbf{riflettono} le rimanenti
        \item Il colore percepito = la $\lambda$ riflessa
        \item Violetto: $\lambda \approx 380$ nm — Rosso: $\lambda \approx 780$ nm
        \item I fotorecettori della retina (coni) rilevano tre bande: rosso, verde, blu
    \end{itemize}
    
\end{frame}
